\section{Antecedents, distinctions and quantitative scope}\label{sec:literature}
\subsection{What the transfer theorem adds}
Conditional projection, small-ball quantization and finite-center approximation are classical; see Graf and Luschgy~\cite{graf}. Predictive representations and future equivalence are developed by Littman, Sutton and Singh~\cite{psr} and Shalizi and Crutchfield~\cite{shalizi}; measurable kernel sufficiency has a categorical treatment in Fritz~\cite{fritz}. Blackwell~\cite{blackwell}, Le Cam~\cite{lecam} and Torgersen~\cite{torgersen} provide the statistical-comparison background. Proposition~\ref{prop:quotient} makes explicit a measurable-image certificate and common report versions; it does not originate predictive sufficiency. The classical total-variation inequality is not the new part of Theorem~\ref{thm:morphism}.

Average contraction of random maps is also classical~\cite{diaconis}. The constructive part of Theorem~\ref{thm:transfer} uses a stronger, specifically physical-law certificate: a cover-preserved relation must control every adapted candidate perturbation across a block. The positive-cone lower bound in Lemma~\ref{lem:filtercover} shows how such a relation can be manufactured for primitive sparse products. Its role is to preserve accumulated positivity while rounding after each raw call. It is sufficient, not an intrinsic necessary condition. Proposition~\ref{prop:repair} allows flagged implementation failures, but does not show that arbitrary local covers or ordinary inwardness imply robust block drift.

The opposite direction now uses actual distributions of continuation-response functions under a joint-law domination certificate. Unlike exact tag separation, it allows a continuum of noisy response functions and charges distortion in the same terminal squared task. Unlike a last-marginal small-ball estimate, it can distinguish a large information cut from an atomic terminal answer. The proof is a Hilbert-space quantization argument; the random-access example explicitly uses an established entropy mechanism~\cite{nayak}. The companion's earlier exact-tag and soft-information arguments remain complementary, rather than being renamed as new theorems.

The optimized acquisition result uses a different constructive mechanism: exact refinement of acquired prefixes, with no rounding of an unretained real posterior. Quantization on Moran sets and measures has an extensive independent theory; for example Zhu~\cite{zhu-moran} proves comparability of optimal Voronoi-cell errors under doubling and separation assumptions. Lemma~\ref{lem:greedy} is a self-contained nonasymptotic product estimate. Its use in Theorems~\ref{thm:adaptive} and \ref{thm:joint} is to identify one scale which also controls every adaptive allocation of paid digit queries. We do not claim a general new quantization theorem for all Moran measures, nonproduct sources or arbitrary adaptive partitions.

\subsection{Decoder side information and functional source coding}
In the Wyner--Ziv problem~\cite{wynerziv}, the encoder observes a memoryless source block $X^n$, while the decoder additionally observes $Y^n$; the rate--distortion expression minimizes $I(X;W)-I(Y;W)$ over a Markov auxiliary $W-X-Y$ and a decoder reproduction satisfying the distortion constraint. At our cut, $H$ is the encoder observation, $U$ is decoder side information, and $z(H,U)$ is the remote conditional-mean task. These identifications are substantive: fixing a causal cut produces an ordinary side-information coding problem. Formula~\eqref{eq:exactfunctional} is its direct finite-code variational form, not a new source-coding principle. Unlike an asymptotic rate, $M$ bounds the entire retained tuple in one actual experiment, and the task and its Bayes baseline remain fixed through the cut.

Orlitsky and Roche~\cite{orlitsky} study reliable computation of a function of the sender's and receiver's combined data. Their problem is closer than ordinary reconstruction of the sender's source: future reports change the function that the earlier information must support. Our response function $z(h,\cdot)$ records exactly that dependence. The difference is not the observation that functions can be compressed. It is the quantitative use of an actual product sublaw to put all responses in a common $L^2(\eta)$ metric, together with retained submeasure mass and recursively attainable bounds under the stated raw acquisition interface. Zero-error distinguishability and asymptotic functional coding are not being asserted equivalent to squared-distortion Hilbert quantization.

Nor is finite blocklength, by itself, a novelty claim. Watanabe, Kuzuoka and Tan~\cite{watanabe} give nonasymptotic achievability and second-order bounds for side-information problems, including Wyner--Ziv coding. Those bounds concern coding probabilities and rates under their information-spectrum formulation. Here \eqref{eq:cutsandwich} instead compares a hard one-cut cardinality directly with a finite-measure geometric distortion, and \eqref{eq:conditionalcut} localizes when the global reference product fails. Lemma~\ref{lem:minorization} derives the reference from physical report kernels. It does not improve the general finite-block coding bounds of that literature.

The recursively executable conclusion is separate from all these one-cut reductions. Theorem~\ref{thm:chart} explicitly stores quantized incoming coordinates, proves an all-tail moment bound, and gives equivalence of online risk and continuation width on its declared correlated sequential class. Theorem~\ref{thm:noisy} computes a positive-noise, continuous-query example in which the same task has different online and checkpoint exponents. Gaussian conditioning, companding and density quantization are standard tools. The claimed synthesis is the raw-kernel continuation lower together with the legal intermediate-state upper and their common-task separation; no independent priority certification is inferred from this comparison.

\subsection{Allocation and successive refinement}
Incremental allocation for a separable discrete resource constraint has classical marginal-analysis antecedents, including Fox~\cite{fox}. Our merged-list argument is an elementary exchange argument in that tradition. In particular, it does not introduce greedy bit allocation for independent components. The objective in \eqref{eq:minmaxscale} is the largest remaining weighted scale, not an unspecified average-rate Lagrangian. Its value also controls the physical-law small balls and the residual uncertainty under every paid progressive policy. These additional links are what make one scale govern both $N$ and $M$ in \eqref{eq:adaptivelaw}.

Equitz and Cover~\cite{equitz} characterize successive refinement by compatible rate--distortion-optimal reproductions, with the coarse reproduction conditionally generated from the finer one. Their staged code describes a known source block at asymptotic rates. Here each new digit must first be physically acquired, and all already acquired digits count against a hard state cardinality rather than an entropy-coded average rate. Exact prefixes provide an executable refinement mechanism, but we do not deduce that arbitrary sources are successively refinable. Nonhomogeneous coordinate scales and singular factors remove a stationary high-rate formula, not the independence assumption. Observed values cannot change future scale values in this model; an extension to genuinely value-dependent experimental design requires another acquisition converse.

\subsection{Operationally close work}
Recursive nonlinear-filter quantization predates this work. Pag\`es and Pham~\cite{pages} and Sellami~\cite{sellami} transport quantization error through filtering recursions. Their approximation constructions do not by themselves identify a fixed cardinality for every retained predictive label or a matching actual-posterior converse. Here Lemmas~\ref{lem:filtercover} and \ref{lem:filterblock} provide that count, and Theorem~\ref{thm:brier} gives the ordinary filtering excess. No novelty is claimed for the Bayes formula, softmax derivative or projective nonexpansiveness separately.

Theorem~9 of Subramanian, Sinha, Seraj and Mahajan~\cite{ais} bounds approximate dynamic-program values from an approximate information-state generator and reward/transition discrepancies. Its policy-loss conclusion concerns control optimization. Theorem~16 of Kara and Y\"uksel~\cite{kara} derives finite-window control bounds under their stability, finite-alphabet and discount hypotheses. Our same-controller physical-law estimate does not imply either optimal-control conclusion. Conversely their conditions do not supply the actual acquired-law lower bound or the hard predictive-state alphabet of \eqref{eq:brier}.

Reachable-belief covering is not a new complexity viewpoint. Zhang, Littman and Chen~\cite{zhang2012} connect covering numbers to POMDP planning and no-reset learning. Theorem~1 of Zhang, Hsu and Lee~\cite{zhang2014} gives a heuristic-guided planning-time upper of order $h\,C(\delta/2)^2$, with depth and resolution chosen from the discount and target value error. Its cover concerns a heuristically reachable search space, not the probability mass actually acquired under a fixed experiment. Grover and Dimitrakakis~\cite{grover} provide adaptive belief discretization and explicit planner-memory/value-error bounds. They therefore cannot be dismissed as not addressing memory; the distinction is planning storage and value approximation versus per-cut predictive alphabet, acquisition uncertainty and a matching geometric lower.

Recent work likewise uses belief geometry operationally. Zhu and Lu~\cite{zhulu} develop offline POMDP evaluation bounds through belief-space coverage and policy stability. This is a learned off-policy evaluation problem, whereas \eqref{eq:central} fixes known kernels and \eqref{eq:adaptivelaw} optimizes a specified progressive sensor. Anjarlekar, Etesami and Srikant~\cite{anjarlekar} use finite-history superstates and policy-based learning with approximation controlled by history length. We do not obtain their learning guarantees from our known-kernel digital implementation. These comparisons are to their published conference versions, not predictions of future results.

Finite-state controller synthesis also has an independent objective. Andriushchenko, \v{C}e\v{s}ka, Junges and Katoen~\cite{inductive} synthesize controllers for POMDP specifications. Hud\'ak and collaborators~\cite{hudak} extract finite-state controllers from learned recurrent policies and evaluate them, including hidden-model robustness. These are substantial finite-controller results, not instances of a terminal codebook. Our converse fixes a common statistical task and the information crossing a cut; it does not establish controller synthesis completeness, robustness for unknown dynamics, or superiority in computational time.

Zero-delay coding and nonanticipative rate-distortion theory address related causal constraints. Wood, Linder and Y\"uksel~\cite{wood}, Ghomi, Linder and Y\"uksel~\cite{ghomi}, and Charalambous, Stavrou and Ahmed~\cite{nonanticipative} study structural or information-constrained formulations. A mutual-information budget is not automatically a hard retained alphabet; nor does our finite-horizon state-order result prove existence of optimal stationary average-cost codes.

\subsection{Scope of the resulting laws}
The new obstruction and the sufficient construction deliberately have different assumptions. Global joint-law domination may have very small mass or fail entirely. Conditional cuts can remain useful, as Example~\ref{ex:diagonal} proves. Theorem~\ref{thm:chart} proves quantitative completeness on its explicit sequential-chart class; passing every available cut test is not proved sufficient outside such a constructively verified class. The profile \eqref{eq:adaptivelaw} is optimized over every adaptive allocation in its raw sensing class, but not over arbitrary interventions, unknown-kernel learning or unbounded stopping. It includes only the stated deterministic phase convention; total persistent state includes any separately stored clock.

For the Gaussian realization, constants deteriorate with $1-\sqrt\kappa$, $1-a$, primitive length, positivity floors, acquired mass and the inverse determinant in \eqref{eq:gaussparams}. For the multiscale realization they depend on $a,d$, not on a nonexistent assumed limiting dimension. Theorems~\ref{thm:digital} and \ref{thm:morphism} charge workspace, program, simulator and time, but do not prove the full optimal feasible region for those coordinates. A lower bound on an output grid is not a lower bound on each internal arithmetic operation. A calibration or deficiency \emph{allowance} alone is not a floor; \eqref{eq:jointlaw} attains such terms only in its explicit inaccessible-sign and erasure family.

Thus the central connection is mathematical rather than terminological: the prepared kernels determine future tests, their measurable predictive realization and actual acquired laws; compatible finite recursions construct the upper; dominated continuation cuts and acquisition variance supply converses; and typed causal morphisms transfer both at their declared resources. The class equivalence and the noisy separation are proved without an independence or contraction assumption on acquired prefix coordinates, but do not constitute a necessary-and-sufficient classification of every recursive acquired geometry.

\paragraph{Serial acquired exponents.}
Theorem~\ref{thm:serial} does not claim new small-ball quantization, Minkowski inequalities or linear error propagation. Its conclusion concerns the same scored experiment at every retained-state cut: observable acquired submeasures supply a lower of each exponent, while global inward covers and a candidate Lyapunov envelope supply a single causal construction. This establishes a conditional compatibility criterion and an explicit ratio of online to checkpoint excess. Neither an asymptotic decoder-side-information rate nor a separate optimal quantizer at each time alone supplies that recursion. The rank-sensitive example verifies the entire sequence of cuts, not only a terminal covariance. In the singular example the preparation geometry and the later task geometry have different exponents, so the maximum must include the final cut. The scope of the criterion is the stated raw-verifiable certificate class; it is not an unconditional duality theorem for all experiments.
